\documentclass[17pt,a4paper,landscape]{extarticle}
\usepackage[margin=2.35cm,top=0.12cm,left=1.05cm]{geometry}
\usepackage{amsmath}
\usepackage{amssymb}
\usepackage{tasks}
\usepackage{tikz}
\usepackage{xcolor}
\usepackage[sfdefault,lf]{FiraSans}
\renewcommand{\familydefault}{\sfdefault}
\setlength{\parindent}{0pt}
\pagestyle{empty}
\settasks{label=\textbf{\arabic*.}, label-width=2.2em, item-indent=2.95em, column-sep=2.1cm, after-item-skip=7.8em}
\renewcommand{\arraystretch}{1.15}
\everymath{\displaystyle}

\begin{document}
\sffamily
\boldmath

\boldmath
{\LARGE \textbf{Writing a statistical paragraph --- Proficient}}\par\vspace{0.8em}
\noindent\textbf{Data set A: Preferred PE activity} Year 9: Basketball 11, Football 8, Netball 7, Volleyball 4. Year 10: Basketball 6, Football 12, Netball 9, Volleyball 5.\\[0.6em]
\textbf{1.} Write a 3-sentence paragraph comparing the two year levels. Include one similarity and one difference.\\[1.2em]
\textbf{Data set B: Daily screen time (students)} 0--1 hr: 3, 1--2 hr: 7, 2--3 hr: 12, 3--4 hr: 8, 4--5 hr: 2.\\[0.6em]
\textbf{2.} Write a 3-sentence paragraph about where most of the data sits. Use exact values.\\[1.2em]
\textbf{Data set C: How students get news (\%)} Social media 45, Family 20, TV 15, Friends 20.\\[0.6em]
\textbf{3.} Write a 3-sentence paragraph that mentions the largest category and one smaller category.\\[1.2em]
\textbf{Data set D: Books borrowed by genre} Mystery 18, Fantasy 14, Comics 12, Science 10, Sport 6.\\[0.6em]
\textbf{4.} Write a paragraph with at least two exact values from data set D.\\[1.2em]
\textbf{Data set E: Homework completed} Class A: Mon 14, Tue 16, Wed 12, Thu 10. Class B: Mon 12, Tue 13, Wed 15, Thu 14.\\[0.6em]
\textbf{5.} Write a 3-sentence comparison paragraph about the two classes.\\[1.2em]
\textbf{Data set F: Bus trip times (minutes)} 0--9: 2, 10--19: 8, 20--29: 11, 30--39: 6, 40--49: 3.\\[0.6em]
\textbf{6.} Write a paragraph and suggest what the data might mean for most students.\\[1.2em]
\textbf{7.} Improve this evidence sentence so it is precise: \emph{Most people chose social media. It is bigger.}\\[1.2em]
\textbf{8.} Explain why a good statistical paragraph needs exact values, not just words like \emph{bigger} or \emph{more}.
\end{document}
